\chapter{Implementação em MATLAB}
\label{ch:matlab}

Neste subcapítulo, são apresentados os resultados obtidos após a junção do cenário de simulação do CoppeliaSim e a implementação da lógica da cinemática direta e inversa do manipulador robótico MyCobot280 em MATLAB. 

\section{Comunicação entre o MATLAB e o CoppeliaSim}
\label{sec:comunicacao}

A simulação corre em modo síncrono (\texttt{sim.setStepping(true)}): só avança um passo, de cerca de 50~ms, quando o MATLAB o pede. Assim, enquanto o MATLAB calcula, o simulador fica parado, e os resultados não dependem da velocidade do computador. A Tabela~\ref{tab:api} mostra as funções usadas no movimento do braço, das classes \texttt{mycobot\_interface} e \texttt{simulator\_interface} fornecidas.

\begin{table}[H]
  \centering
  \caption{Funções usadas para mover o braço e chamadas à API remota correspondentes.}
  \label{tab:api}
  \small
  \begin{tabular}{p{4.4cm}p{4.0cm}p{3.6cm}}
    \hline
    Função MATLAB & Chamada à API & Para quê \\[2pt]
    \hline
    \texttt{robot\_arm.set\_joints(q)} & \texttt{sim.setJoint}\allowbreak\texttt{TargetPosition} & enviar os ângulos pretendidos \\
    \texttt{robot\_arm.get\_joints()} & \texttt{sim.getJointPosition} & ler os ângulos das juntas \\
    \texttt{sim\_interface.}\allowbreak\texttt{step\_simulation()} & \texttt{sim.step} & avançar um passo \\
    \texttt{sim\_interface.}\allowbreak\texttt{get\_simulation\_time()} & \texttt{sim.getSimulationTime} & tempo de simulação \\
    \hline
  \end{tabular}
\end{table}

As leituras das juntas são o que os \textit{encoders} do braço real dariam. As posições reais dos objetos, que só o simulador conhece, são usadas apenas nos testes de diagnóstico, para medir os erros.

\section{Organização do código}
\label{sec:organizacao}
 
\begin{table}[H]
  \centering
  \caption{Funções do pacote \texttt{+movement} usadas nesta componente.}
  \label{tab:funcoes}
  \begin{tabular}{ll}
    \hline
    Função & O que faz \\
    \hline
    \texttt{direct\_kin} & cinemática direta (Capítulo~\ref{ch:dir_kin}) \\
    \texttt{ik\_2dof} & braço planar de 2 DOF, como na aula \\
    \texttt{inverse\_kin} & cinemática inversa (Capítulo~\ref{ch:inv_kin}) \\
    \texttt{move\_joints} & movimento suave no espaço das juntas \\
    \texttt{move\_linear} & movimento da garra em linha reta \\
    \texttt{go\_home} & regresso à pose inicial a partir de qualquer pose \\
    \texttt{plan\_vertical} & escolha de uma inclinação para toda uma descida \\
    \hline
  \end{tabular}
\end{table}

\section{Discussão}
\label{sec:disc}
 
\subsection{Validação do modelo}
 
O erro do modelo compara a posição calculada pela cinemática direta, com as juntas lidas no simulador, com a posição real do centro de preensão. Ficou entre 1,2 e 4,0~mm nos alvos testados, abaixo da tolerância de 5~mm, o que valida em conjunto a função \texttt{direct\_kin}, a conversão \texttt{world\_to\_base} e os comprimentos dos elos medidos no cenário. As diferenças que restam são pequenas e podem vir de os comprimentos medidos não coincidirem exatamente com a geometria do modelo do robô no CoppeliaSim e de o centro de preensão ser definido a partir de um ponto do modelo. Este erro não depende do movimento, pois só compara dois cálculos da mesma pose, e por isso mede a qualidade do modelo e não a do controlo.
 
\subsection{Erro de alcance e erro nas juntas}
 
O erro de alcance compara o ponto pedido com a posição real da garra, pelo que mistura dois efeitos: o erro do modelo e o facto de as juntas não chegarem exatamente aos ângulos pedidos. O erro nas juntas isola o segundo. Se o erro de alcance for semelhante ao erro do modelo e o erro nas juntas for pequeno, o movimento atinge o ponto pedido e o erro restante vem do modelo. Se o erro de alcance for maior, a diferença vem de as juntas ficarem ainda a caminho ou abaixo do ângulo pedido, por exemplo por efeito da gravidade ou do tempo de estabilização, de 2~s, ser curto para o movimento em causa.
\subsection{Alcance e inclinação da garra}
 
A cinemática inversa tenta primeiro a inclinação preferida, com a garra na vertical, e só se esta falhar inclina a garra. Com a implementação que usámos, para alvos a poucos centímetros da mesa, o alcance horizontal é de cerca de 27--28~cm do eixo da base com a garra na vertical e passa a cerca de 36--37~cm quando se permite inclinar a garra até $30^\circ$ da horizontal. Os alvos mais afastados da base são, por isso, os que exigem inclinar a garra, e são também os que deixam o braço quase esticado, onde pequenos erros nas juntas se traduzem em erros maiores de posição. Por outro lado, um alvo demasiado perto do eixo da base pode ser inalcançável, porque o desvio lateral $L_5$ impede a garra de chegar ao eixo de J1.
 
\subsection{Limitações}
 
O movimento \texttt{move\_joints} interpola as juntas, e não a posição da garra, pelo que a trajetória do \textit{end-effector} entre dois pontos não é uma reta; passa pela pose inicial para evitar a mesa, mas não é uma trajetória planeada no espaço cartesiano. Quando é preciso um movimento em linha reta, usa-se \texttt{move\_linear}, que resolve a cinemática inversa a cada passo. Além disso, os resultados são de simulação e usam as posições reais dos objetos, que só o simulador conhece, apenas para definir os alvos e medir os erros; no braço real, os alvos viriam da câmara e não existiria o erro do modelo medido desta forma.
 

\section{Cinemática direta}
\label{sec:cin_dir}

A função \texttt{direct\_kin} serve para implementar a Equação~\eqref{eq:dk_final}, seguindo a mesma ordem de dedução utilizada anteriormente.

\begin{lstlisting}[language=Matlab, caption={Função \texttt{direct\_kin}}, label={lst:dir_kin}]
function [tip, approach] = direct_kin(q, Links)
    q = q(:)';
    if numel(q) == 4
        q = [q 0 0];                % bloquear J5 e J6 
    end
    L = Links;
    D = L(6) + L(7);    

    thetaA = q(2) + pi/2;           % elo L2 (J2 -> J3)
    thetaB = thetaA + q(3);         % elo L3 (J3 -> J4)
    thetaC = thetaB + q(4);         % elo L4 (J4 -> J5)
    alpha  = thetaC - pi/2;         % garra: perpendicular a L4 (-90 graus = para baixo)

    % J2 (ombro)
    d = 0;                          h = L(1);
    % J3 (cotovelo)
    d = d + L(2)*cos(thetaA);       h = h + L(2)*sin(thetaA);
    % J4 (pulso)
    d = d + L(3)*cos(thetaB);       h = h + L(3)*sin(thetaB);
    % J5
    d = d + L(4)*cos(thetaC);       h = h + L(4)*sin(thetaC);
    % centro de preensao
    d = d + D*cos(q(5))*cos(alpha); h = h + D*cos(q(5))*sin(alpha);
    l = L(5) - D*sin(q(5));        

    x = l*cos(q(1)) + d*cos(q(1) + pi/2);  % = l*cos(q1) - d*sin(q1)
    y = l*sin(q(1)) + d*sin(q(1) + pi/2);  % = l*sin(q1) + d*cos(q1)
    z = h;
    tip = [x; y; z];

    aL = -sin(q(5));                % componente lateral
    aD =  cos(q(5))*cos(alpha);     % componente em frente
    aH =  cos(q(5))*sin(alpha);     % componente para cima
    approach = [aL*cos(q(1)) + aD*cos(q(1) + pi/2);
                aL*sin(q(1)) + aD*sin(q(1) + pi/2);
                aH];
end
\end{lstlisting}

\begin{table}[H]
  \centering
  \caption{Entradas e saídas de \texttt{direct\_kin}.}
  \label{tab:direct_io}
  \begin{tabular}{ll}
    \toprule
    Argumento & Significado \\
    \midrule
    \multicolumn{2}{l}{\textit{Entradas}} \\
    \texttt{q}     & ângulos das juntas [rad], $1\times4$ ou $1\times6$ \\
    \texttt{Links} & os sete comprimentos da Tabela~\ref{tab:elos} [m] \\
    \midrule
    \multicolumn{2}{l}{\textit{Saídas}} \\
    \texttt{tip}      & posição $3\times1$ do centro de preensão, referencial da base [m] \\
    \texttt{approach} & vetor unitário da direção da garra \\
    \bottomrule
  \end{tabular}
\end{table}

Esta função é utilizada em vários pontos da execução para identificar a posição da garra a partir das leituras das juntas

\section{Cinemática inversa}
\label{sec:cin_inv}

A função \texttt{inverse\_kin} implementa a formulação matemática da cinemática inversa apresentada no Capítulo~\ref{ch:inv_kin}, permitindo determinar os ângulos articulares necessários para alcançar uma posição pretendida do centro de preensão.

A resolução é realizada com o auxílio das funções \texttt{ik\_2dof} e \texttt{solve\_pitch}. A primeira resolve o problema geométrico de um manipulador planar de dois graus de liberdade, enquanto a segunda determina os ângulos das juntas J1 a J4 para uma inclinação específica da garra, verificando as restrições de alcance e os limites articulares.

A função principal começa por calcular a distância frontal e a altura relativamente ao ombro. Seguidamente, procura uma inclinação válida para a garra, começando pela orientação preferencial de $90^\circ$ e testando, quando necessário, inclinações alternativas em incrementos de $5^\circ$.

Quando é encontrada uma solução válida numa inclinação diferente da preferencial, é realizado um refinamento através do método da bisseção. Por fim, é calculada a rotação da junta J6, caso tenha sido fornecida a orientação do objeto e a inclinação da garra satisfaça a condição estabelecida.

A implementação da função principal encontra-se apresentada na Listagem~\ref{lst:inv_kin}.

\begin{lstlisting}[
    language=Matlab,
    caption={Implementação da função inverse\_kin.},
    label={lst:inv_kin}
]
function [armJoints, isReachable, pitchUsed] = ...
    inverse_kin(x, y, z, Links, qMin, qMax, ...
                objYaw, prefPitch, family, fixedPitch)

    % Valores por defeito
    if nargin < 5 || isempty(qMin)
        qMin = deg2rad([-165 -165 -165 -165 -165 -175]);
    end
    if nargin < 6 || isempty(qMax)
        qMax = deg2rad([165 165 165 165 165 175]);
    end
    if nargin < 7
        objYaw = [];
    end
    if nargin < 8 || isempty(prefPitch)
        prefPitch = pi/2;
    end
    if nargin < 9 || isempty(family)
        family = 'up';
    end
    if nargin < 10 || isempty(fixedPitch)
        fixedPitch = false;
    end

    L = Links;
    D = L(6) + L(7);

    armJoints = zeros(1,6);
    isReachable = false;
    pitchUsed = NaN;

    % Plano horizontal
    r = hypot(x,y);
    if r <= L(5)
        return;
    end

    dFwd = sqrt(r^2-L(5)^2);
    hJ2 = z-L(1);

    % Restricoes
    ELBOW_MAX = deg2rad(140);
    PITCH_MAX = deg2rad(135);
    PITCH_MIN = deg2rad(30);
    STEP = deg2rad(5);

    if ~strcmp(family,'down')
        family = 'up';
    end

    if fixedPitch
        dList = 0;
    else
        dList = 0:STEP:(PITCH_MAX-PITCH_MIN);
    end

    % Procurar inclinacao valida
    for d = dList
        if d == 0
            cands = prefPitch;
        else
            cands = [prefPitch+d, prefPitch-d];
        end

        for psi = cands
            if ~fixedPitch && ...
                    (psi > PITCH_MAX || psi < PITCH_MIN)
                continue;
            end

            best = solve_pitch(psi,x,y,dFwd,hJ2, ...
                L,D,qMin,qMax,ELBOW_MAX,family);

            if isempty(best)
                continue;
            end

            % Refinamento por bissecao
            if d > 0
                good = psi;
                bad = psi-sign(psi-prefPitch)*STEP;

                for it = 1:15
                    mid = (good+bad)/2;
                    qm = solve_pitch(mid,x,y,dFwd,hJ2, ...
                        L,D,qMin,qMax,ELBOW_MAX,family);

                    if isempty(qm)
                        bad = mid;
                    else
                        good = mid;
                        best = qm;
                    end
                end
                psi = good;
            end

            % Orientacao da junta J6
            q6 = 0;
            if ~isempty(objYaw) && psi > deg2rad(45)
                q6 = mod(best(1)-objYaw+pi/4,pi/2)-pi/4;
            end

            armJoints = [best 0 q6];
            isReachable = true;
            pitchUsed = psi;
            return;
        end
    end
end
\end{lstlisting}

A função devolve os seis ângulos articulares, um indicador de existência de solução e a inclinação efetivamente utilizada. Quando não é encontrada uma configuração válida, o indicador \texttt{isReachable} permanece falso.

O argumento \texttt{fixedPitch} permite desativar a procura automática de inclinações alternativas, sendo utilizado, por exemplo, em movimentos verticais nos quais se pretende manter a mesma orientação da garra ao longo da trajetória.

As entradas e saídas da função encontram-se resumidas na Tabela~\ref{tab:inverse_io}.

\begin{table}[H]
    \centering
    \caption{Entradas e saídas da função \texttt{inverse\_kin}.}
    \label{tab:inverse_io}
    \small
    \begin{tabular}{lp{8.5cm}}
        \hline
        Argumento & Significado \\
        \hline
        \multicolumn{2}{l}{\textit{Entradas}} \\
        \texttt{x, y, z} & Posição pretendida no referencial da base [m] \\
        \texttt{Links} & Parâmetros geométricos do manipulador [m] \\
        \texttt{qMin, qMax} & Limites articulares [rad] \\
        \texttt{objYaw} & Orientação do objeto em torno de $z$, opcional [rad] \\
        \texttt{prefPitch} & Inclinação preferencial da garra, opcional [rad] \\
        \texttt{family} & Configuração \texttt{'up'} ou \texttt{'down'}, opcional \\
        \texttt{fixedPitch} & Utilização exclusiva da inclinação indicada, opcional \\
        \hline
        \multicolumn{2}{l}{\textit{Saídas}} \\
        \texttt{armJoints} & Vetor $[q_1\ q_2\ q_3\ q_4\ q_5\ q_6]$ [rad] \\
        \texttt{isReachable} & Indicador de existência de solução válida \\
        \texttt{pitchUsed} & Inclinação efetivamente utilizada [rad] \\
        \hline
    \end{tabular}
\end{table}

\subsection{Funções auxiliares da cinemática inversa}

A função \texttt{ik\_2dof} resolve a cinemática inversa do braço planar formado pelos elos $L_2$ e $L_3$, através da lei dos cossenos. A implementação permite selecionar a configuração com o cotovelo para cima ou para baixo, conforme descrito na Secção~\ref{sec:ik_lateral}.

Por sua vez, a função \texttt{solve\_pitch} recebe uma inclinação $\psi$ previamente definida, calcula a posição do pulso J4 e utiliza \texttt{ik\_2dof} para determinar os ângulos das juntas J2 e J3. Seguidamente, determina J1 e J4 e verifica se a configuração obtida respeita os limites articulares, a flexão máxima do cotovelo e as restrições de altura das juntas.

As implementações destas funções encontram-se nas Listagens~\ref{lst:ik2dof} e~\ref{lst:solvepitch}.

\begin{lstlisting}[
    language=Matlab,
    caption={Implementação da função ik\_2dof.},
    label={lst:ik2dof}
]
function [theta1, theta2, isReachable] = ...
    ik_2dof(xd, yd, LA, LB, solution_type)

    theta1 = NaN;
    theta2 = NaN;
    isReachable = false;

    r = sqrt(xd^2+yd^2);
    alpha1 = atan2(yd,xd);

    arg1 = (LA^2+r^2-LB^2)/(2*LA*r);
    arg2 = (LA^2+LB^2-r^2)/(2*LA*LB);

    if r == 0 || abs(arg1) > 1 || abs(arg2) > 1
        return;
    end

    beta1 = acos(arg1);
    beta2 = acos(arg2);

    if strcmp(solution_type,'down')
        theta1 = alpha1-beta1;
        theta2 = pi-beta2;
    else
        theta1 = alpha1+beta1;
        theta2 = -(pi-beta2);
    end

    isReachable = true;
end
\end{lstlisting}

\begin{lstlisting}[
    language=Matlab,
    caption={Implementação da função solve\_pitch.},
    label={lst:solvepitch}
]
function q = solve_pitch(psi,x,y,dFwd,hJ2, ...
    L,D,qMin,qMax,elbowMax,family)

    wrap = @(t) atan2(sin(t),cos(t));
    q = [];

    % Angulo da base
    q1 = atan2(y,x)-atan2(dFwd,L(5));

    % Posicao do pulso
    alpha3 = -psi;
    x3d = dFwd-D*cos(alpha3) ...
        -L(4)*cos(alpha3+pi/2);
    y3d = hJ2-D*sin(alpha3) ...
        -L(4)*sin(alpha3+pi/2);

    % Braco planar de dois elos
    [theta1,theta2,ok] = movement.ik_2dof( ...
        x3d,y3d,L(2),L(3),family);

    if ~ok
        return;
    end

    % Angulo da junta J4
    theta3 = (alpha3+pi/2)-theta1-theta2;

    % Angulos articulares
    q = wrap([q1,theta1-pi/2,theta2,theta3]);

    % Alturas das juntas
    zElbow = L(1)+L(2)*cos(q(2));
    zJ4 = zElbow+L(3)*cos(q(2)+q(3));
    zJ5 = zJ4+L(4)*cos(q(2)+q(3)+q(4));

    % Validacao da configuracao
    if any(q < qMin(1:4)) || ...
            any(q > qMax(1:4)) || ...
            abs(q(3)) > elbowMax || ...
            min([zElbow zJ4 zJ5]) <= 0
        q = [];
    end
end
\end{lstlisting}

As entradas e saídas das duas funções auxiliares encontram-se apresentadas nas Tabelas~\ref{tab:ik2dof_io} e~\ref{tab:solvepitch_io}.

\begin{table}[H]
    \centering
    \caption{Entradas e saídas da função \texttt{ik\_2dof}.}
    \label{tab:ik2dof_io}
    \small
    \begin{tabular}{lp{8cm}}
        \toprule
        Argumento & Significado \\
        \midrule
        \multicolumn{2}{l}{\textit{Entradas}} \\
        \texttt{xd, yd} & Coordenadas do pulso relativamente ao ombro [m] \\
        \texttt{LA, LB} & Comprimentos dos dois elos [m] \\
        \texttt{solution\_type} & Configuração \texttt{'up'} ou \texttt{'down'} \\
        \midrule
        \multicolumn{2}{l}{\textit{Saídas}} \\
        \texttt{theta1} & Orientação absoluta do primeiro elo [rad] \\
        \texttt{theta2} & Rotação do segundo elo relativamente ao primeiro [rad] \\
        \texttt{isReachable} & Indicador de solução geométrica válida \\
        \bottomrule
    \end{tabular}
\end{table}

\begin{table}[H]
    \centering
    \caption{Entradas e saídas da função \texttt{solve\_pitch}.}
    \label{tab:solvepitch_io}
    \small
    \begin{tabular}{lp{8cm}}
        \toprule
        Argumento & Significado \\
        \midrule
        \multicolumn{2}{l}{\textit{Entradas}} \\
        \texttt{psi} & Inclinação definida para a garra [rad] \\
        \texttt{x, y} & Coordenadas horizontais do alvo [m] \\
        \texttt{dFwd, hJ2} & Distância frontal e altura relativa ao ombro [m] \\
        \texttt{L, D} & Parâmetros geométricos do manipulador [m] \\
        \texttt{qMin, qMax} & Limites articulares [rad] \\
        \texttt{elbowMax} & Flexão máxima admitida para o cotovelo [rad] \\
        \texttt{family} & Configuração \texttt{'up'} ou \texttt{'down'} \\
        \midrule
        \multicolumn{2}{l}{\textit{Saída}} \\
        \texttt{q} & Vetor $[q_1\ q_2\ q_3\ q_4]$ [rad] ou vazio \\
        \bottomrule
    \end{tabular}
\end{table}

Deste modo, a implementação separa a procura de uma inclinação válida da resolução geométrica para uma inclinação específica, permitindo reutilizar os cálculos cinemáticos em diferentes procedimentos de movimento.

A função \texttt{inverse\_kin} é utilizada tanto no posicionamento do manipulador como no planeamento de trajetórias, sendo chamada repetidamente quando é necessário determinar configurações articulares para diferentes pontos do espaço de trabalho.


\section{Movimento até uma posição alvo}
\label{sec:imp_move}

O movimento do manipulador até uma posição pretendida envolve a transformação das coordenadas do alvo, o cálculo da configuração articular e a execução do movimento através das funções desenvolvidas em MATLAB.

Quando a posição do alvo é fornecida no referencial do mundo do CoppeliaSim, o procedimento pode ser dividido nas seguintes etapas:

\begin{enumerate}
    \item \textbf{Transformação de coordenadas:}
    a posição pretendida $(X_{\mathrm{target}},Y_{\mathrm{target}},Z_{\mathrm{target}})$ é convertida para o referencial da base do manipulador através da função \texttt{aux.world\_to\_base}, considerando a translação e a rotação entre os dois referenciais;

    \item \textbf{Cálculo da configuração articular:}
    a função \texttt{inverse\_kin} determina os ângulos necessários para alcançar a posição pretendida, respeitando as restrições geométricas e articulares definidas. Caso não exista uma solução válida, o movimento pode ser rejeitado antes da sua execução;

    \item \textbf{Execução do movimento:}
    após determinar uma configuração válida, o manipulador pode ser deslocado através da função \texttt{move\_joints}, que realiza a interpolação dos ângulos articulares, ou da função \texttt{move\_linear}, utilizada quando se pretende seguir uma trajetória aproximadamente retilínea no espaço cartesiano;

    \item \textbf{Verificação do posicionamento:}
    os ângulos articulares obtidos durante a execução podem ser utilizados pela função \texttt{direct\_kin} para calcular a posição do centro de preensão e avaliar a diferença relativamente à posição pretendida.
\end{enumerate}

\subsection{Movimento no espaço das juntas}

A função \texttt{move\_joints} permite deslocar o manipulador entre duas configurações articulares, garantindo uma evolução suave dos ângulos das juntas.

Em vez de enviar diretamente os ângulos finais, a função calcula, em cada passo de simulação, uma configuração intermédia entre a configuração inicial e a pretendida. Todas as juntas iniciam e terminam o movimento ao mesmo tempo, independentemente da amplitude da rotação de cada uma.

A evolução entre as duas configurações segue um perfil suave, em forma de S: as juntas arrancam devagar, aceleram até ao meio do movimento e abrandam até parar na posição pretendida. Desta forma, o movimento começa e termina com velocidade nula, evitando alterações bruscas da velocidade e solavancos no manipulador.

A duração do movimento é definida pela junta que tem de rodar mais. Esta duração é calculada de modo a que a velocidade dessa junta nunca ultrapasse a velocidade máxima pretendida, de $30^\circ$/s por defeito. As restantes juntas, como rodam menos no mesmo tempo, movem-se mais devagar.

A execução é realizada de forma sincronizada, sendo enviadas configurações articulares intermédias ao manipulador em cada passo de simulação.

No final, a função verifica se todas as juntas se encontram a menos de $1^\circ$ dos respetivos valores pretendidos.
\subsection{Movimentos em linha reta}

Para movimentos de aproximação a objetos, nomeadamente durante a descida da garra até uma peça, é utilizada a função \texttt{move\_linear}.

Ao contrário do movimento no espaço das juntas, esta função faz evoluir um ponto de referência ao longo da direção que une a posição atual do centro de preensão à posição pretendida.

Em cada passo de simulação, o programa atualiza a velocidade do ponto de referência, respeitando a velocidade máxima definida e a aceleração configurada através do parâmetro \texttt{cfg.ACC}.

A velocidade é aumentada gradualmente e reduzida à medida que o ponto se aproxima do destino, de forma a evitar uma paragem brusca.

Para cada posição intermédia, a função \texttt{inverse\_kin} calcula os ângulos articulares correspondentes, que são enviados ao manipulador.

Após o ponto de referência atingir o destino, o programa compara a posição pretendida com a posição calculada pela cinemática direta a partir dos ângulos efetivamente lidos.

Quando necessário, é aplicada uma correção progressiva ao ponto de referência para compensar diferenças de posicionamento, incluindo os efeitos associados à gravidade.

O movimento é considerado concluído com sucesso quando o erro de posicionamento é inferior a $3$ mm, dentro das condições de execução estabelecidas.

\subsection{Planeamento de movimentos verticais}

Antes de executar determinados movimentos verticais, a função \texttt{plan\_vertical} seleciona uma inclinação fixa da garra que permita percorrer toda a trajetória.

Este procedimento é particularmente importante durante a aproximação a objetos situados perto da base do manipulador, onde pequenas alterações de posição podem exigir variações significativas dos ângulos articulares.

O algoritmo testa diferentes inclinações, começando pela orientação vertical para baixo, correspondente a $\psi=90^\circ$.

Para cada inclinação candidata, são verificadas sucessivamente as posições ao longo da trajetória vertical, em intervalos de $2$ mm.

Uma inclinação apenas é aceite quando são satisfeitas as seguintes condições:

\begin{itemize}
    \item Existe uma solução válida de cinemática inversa para todos os pontos da trajetória;

    \item A variação de cada uma das juntas J1 a J4 entre posições consecutivas não ultrapassa $5^\circ$;

    \item Quando a verificação de colisões está ativada, as configurações articulares não provocam colisões do manipulador consigo próprio.
\end{itemize}

Durante esta verificação, é utilizada a configuração com o cotovelo para cima, sendo mantida a mesma inclinação da garra ao longo de todo o movimento.

Esta abordagem reduz a possibilidade de alterações abruptas das configurações articulares durante a descida.

A Figura~\ref{fig:move} apresenta uma representação esquemática da sequência de aproximação a uma peça.

\begin{figure}[htbp]
    \centering
    % Sequência do movimento até uma peça, vista lateral (esboço, sem escala)
\begin{tikzpicture}[scale=1.2, >=stealth,
    elo/.style={line width=2pt, line cap=round},
    junta/.style={circle, draw, fill=white, line width=0.9pt, inner sep=1.6pt},
    pt/.style={circle, fill, inner sep=1.5pt},
    every node/.style={font=\footnotesize}]

  % mesa e base
  \fill[pattern=north east lines, pattern color=gray!60] (-1.2,-0.25) rectangle (10.2,0);
  \draw[thick] (-1.2,0) -- (10.2,0) node[below left=2pt and 0pt] {mesa};
  \fill[gray!30] (-0.45,0) rectangle (0.45,0.5);
  \draw (-0.45,0) rectangle (0.45,0.5);
  % braço na pose inicial (esboço)
  \coordinate (J2) at (0,1.6);
  \draw[elo, gray!70] (0,0.5) -- (J2) -- (0.3,3.3) -- (1.6,4.1);
  \node[junta] at (J2) {}; \node[junta] at (0.3,3.3) {};
  % peça
  \fill[red!65] (6.3,0) rectangle (7.1,0.8);
  \draw (6.3,0) rectangle (7.1,0.8);

  \coordinate (P0) at (1.6,4.1);
  \coordinate (P1) at (6.7,3.0);
  \coordinate (P2) at (6.7,1.4);

  % 1: move_joints (trajetória curva da garra)
  \draw[->, blue!70!black, line width=1.1pt, dashed] (P0) .. controls (3.6,5.2) and (5.8,4.4) .. (P1);
  \node[blue!70!black, align=center] at (3.9,5.35) {\textbf{1.} \texttt{move\_joints}};
  % 2: move_linear (descida vertical)
  \draw[->, orange!85!black, line width=1.3pt] (P1) -- (P2);
  \foreach \h in {2.8,2.6,...,1.5} \draw[orange!85!black] (6.6,\h) -- (6.8,\h);
  \node[orange!85!black, align=left, anchor=west] at (7.7,2.4) {\textbf{2.} \texttt{move\_linear}\\descida em linha reta,\\uma só inclinação\\(\texttt{plan\_vertical},\\de 2 em 2 mm)};
  % 3: subida e go_home
  \draw[->, green!45!black, line width=1.1pt, dotted] (6.2,1.5) -- (6.2,2.9) .. controls (5.0,4.0) and (3.0,4.6) .. (1.8,4.25);
  \node[green!45!black, align=center] at (3.5,3.1) {\textbf{3.} \texttt{go\_home}: sobe na\\vertical e regressa};

  \node[pt] at (P0) {}; \node[above left=1pt and 0pt] at (P0) {pose inicial};
  \node[pt] at (P1) {}; \node[above right=5pt and 2pt] at (P1) {aproximação};
  \node[pt] at (P2) {};
  \draw[<->, gray!70!black, thin] (7.35,0) -- (7.35,1.4) node[midway, right, black] {2 cm};
  \draw[gray!70!black, thin, dashed] (6.7,1.4) -- (7.45,1.4);
\end{tikzpicture}
    \caption{Representação esquemática da sequência de movimento até uma peça, numa vista lateral e sem escala.}
    \label{fig:move}
\end{figure}

\subsection{Regresso à pose inicial}

A função \texttt{go\_home} permite deslocar o manipulador para uma configuração inicial previamente definida, independentemente da posição articular de partida, desde que o movimento seja realizável.

O procedimento começa pela abertura da garra e pela leitura dos ângulos articulares atuais.

Através da função \texttt{direct\_kin}, é determinada a posição do centro de preensão. Caso a sua altura seja inferior ao valor de segurança definido, de $0{,}15$ m relativamente à base, o algoritmo procura realizar uma subida gradual antes do regresso à configuração inicial.

Durante esta subida, são calculadas sucessivamente novas configurações articulares através da função \texttt{inverse\_kin}, procurando manter a inclinação inicial da garra.

Se não for possível encontrar uma configuração válida ou se forem detetadas variações articulares excessivas, a subida é interrompida e o programa avança para a etapa seguinte.

Posteriormente, a função \texttt{move\_joints} é utilizada para executar o movimento suave até à configuração inicial.

Por fim, os ângulos articulares são novamente lidos, sendo o regresso considerado bem-sucedido quando todas as juntas apresentam um desvio inferior a $1^\circ$ relativamente aos valores pretendidos.

Este procedimento baseia-se nas leituras articulares, permitindo utilizar a mesma lógica de controlo com a interface do manipulador simulado ou com uma interface compatível de um braço físico.
